\chapter{周期固体、基组与赝势}
\sourcepages{原英文稿第 10 章；朋友的 50 页笔记未覆盖本章}
\section{平移对称性如何减少未知量}
晶体有无限多个原胞，逐原子求解不可行。设原胞平移向量为 $\bm a_1,\bm a_2,\bm a_3$，晶格向量为 $\bm R=\sum_i n_i\bm a_i$，有效势满足 $v(\bm r+\bm R)=v(\bm r)$。定义平移算符 $T_R\psi(\bm r)=\psi(\bm r+\bm R)$。势的周期性以及动能的平移不变性给出 $[H,T_R]=0$。对易来自 Hamilton 算符的对称性，不能仅由 $T_R$ 是酉算符推出。

在共同本征态上写 $T_R\psi=t_R\psi$。酉性保证 $|t_R|=1$，群乘法 $T_RT_{R'}=T_{R+R'}$ 要求本征值满足相同乘法关系，故可写 $t_R=e^{\ii\bm k\cdot\bm R}$。定义 $u_{n\bm k}=e^{-\ii\bm k\cdot\bm r}\psi_{n\bm k}$，立即得到
\begin{equation}
\psi_{n\bm k}(\bm r)=e^{\ii\bm k\cdot\bm r}u_{n\bm k}(\bm r),\qquad u_{n\bm k}(\bm r+\bm R)=u_{n\bm k}(\bm r).
\end{equation}
这就是 Bloch 形式，$n$ 标记能带。同一个 Bloch 态的模平方在原胞间相同，但波函数本身一般只差相位，而非严格原胞周期。

倒格矢满足 $\bm a_i\cdot\bm b_j=2\pi\delta_{ij}$。令波函数在 $N_i\bm a_i$ 形成的大超胞上周期，则允许 $\bm k=\sum_i(m_i/N_i)\bm b_i$。相差倒格矢的 $\bm k$ 描述等价平移本征值，故只需第一 Brillouin 区。将 Bloch 形式代入 KS 方程，得到原胞内方程
\begin{equation}
\left[-\tfrac12(\nabla+\ii\bm k)^2+v_s(\bm r)\right]u_{n\bm k}=\varepsilon_{n\bm k}u_{n\bm k}.
\end{equation}
无限晶体问题转化为不同 $\bm k$ 参数下的有限原胞问题。

\section{积分采样与能带路径不是同一件事}
每原胞的 Brillouin 区平均为 $\Omega_{BZ}^{-1}\int_{BZ}A(\bm k)\dd^3k$，其中 $\Omega_{BZ}=(2\pi)^3/\Omega_{\rm cell}$。数值上用 $\sum_k w_kA(\bm k)$，权重和为一。若 $u$ 在原胞归一化，电子密度为 $n(\bm r)=\sum_{nk}w_k f_{nk}|u_{nk}(\bm r)|^2$，自旋简并是否计入 $f$ 必须前后一致。

Monkhorst--Pack 网格的一种写法是 $\bm k=\sum_i(2m_i-N_i-1)\bm b_i/(2N_i)$，$m_i=1,\ldots,N_i$。晶体对称性可把等价点合并，但合并权重必须来自整个离散网格。画能带的高对称路径只是一条线，不能代替三维积分网格。金属的占据在 Fermi 面附近变化快，通常比绝缘体需要更细的采样与展宽收敛检查。

无磁、无自旋轨道耦合且势为实标量时，复共轭把 $\bm k$ 态映为同能量的 $-\bm k$ 态。在合适基中 $\Gamma$ 点可选实轨道。这不意味着所有具有时间反演对称性的系统都有 Dirac 锥；能带交叉还依赖其他对称性和能带结构。含自旋时，时间反演也不能简单等同于复共轭。

\section{平面波为什么方便}
周期函数可展开为 $u_{n\bm k}=\Omega_{\rm cell}^{-1/2}\sum_Gc_{n\bm k}(G)e^{\ii\bm G\cdot\bm r}$。由于 $\nabla e^{\ii G r}=\ii\bm G e^{\ii G r}$、$\nabla^2e^{\ii G r}=-G^2e^{\ii G r}$，动能在倒空间对角。保留 $|\bm k+\bm G|^2/2\le E_{\rm cut}$ 的项，把无限基变成有限基。增大截断在固定几何、嵌套基和充分求解下扩大变分空间，但实际计算还要控制积分网格等误差。

势能乘法在实空间方便，故常以 FFT 在两种表示之间切换。密度包含 $c^*(G)c(G')e^{\ii(G'-G)r}$；若轨道波矢最大为 $G_{\max}$，密度包含至 $2G_{\max}$ 的成分，对应四倍动能截断。这是平面波乘积的几何结论；超软赝势或 PAW 增广密度可能需要更高密度截断，不能机械设成四倍。

Poisson 方程 $\nabla^2v_H=-4\pi n$ 给出 $v_H(G)=4\pi n(G)/G^2$，仅对 $G\ne0$ 成立。电子密度的 $n(0)$ 是平均电子数密度，\emph{不为零}。周期静电的零模需要结合离子电荷、中性条件和势零点约定处理。带电超胞还涉及补偿背景，不能把零模发散解释成分子也变为零。

\section{赝势为何与基组选择相连}
核附近全电子轨道快速变化，并与芯态正交；直接用平面波表示需很高截断。赝势在芯区内用平滑价态与有效作用替代显式芯电子，同时希望保持芯区外的散射性质。对某个角动量通道，模守恒构造要求参考能量和芯区外径向函数匹配，并满足
\begin{equation}
\int_0^{r_c}|R_l(r)|^2r^2\dd r=\int_0^{r_c}|\widetilde R_l(r)|^2r^2\dd r.
\end{equation}
这是约束之一，不是“任意平滑后就等价”的保证。不同角动量需要不同作用，因而赝势一般包含非局域投影项。可转移性必须在不同价态、成键和压强环境中检验。

局域原子轨道基则以较少函数描述成键，但需要控制基组完备性与近线性相关。Slater 型径向函数有 $e^{-\zeta r}$ 的衰减，Gaussian 型用多项式乘 $e^{-\alpha r^2}$，便于积分运算。移动原子会移动局域基，所以力中要考虑第四章的 Pulay 项。实空间有限差分是另一种选择，例如 $\partial_x^2\psi\simeq[\psi(x+a)-2\psi(x)+\psi(x-a)]/a^2$；其网格误差需单独收敛，通常不具有相同的严格变分单调性。

\section{从能带曲率理解有效质量}
在波包半经典近似下，群速度为 $v_\alpha=\hbar^{-1}\partial\varepsilon/\partial k_\alpha$，外力满足 $\hbar\dot k_\beta=F_\beta$。对速度再求时间导数，得到
\begin{equation}
\dot v_\alpha=\sum_\beta\frac1{\hbar^2}\frac{\partial^2\varepsilon}{\partial k_\alpha\partial k_\beta}F_\beta,
\qquad (m^{*-1})_{\alpha\beta}=\frac1{\hbar^2}\frac{\partial^2\varepsilon}{\partial k_\alpha\partial k_\beta}.
\end{equation}
有效质量是带边曲率张量，不是先验输入来决定能带。仅在各向同性抛物线附近，才可简化成一个标量 $m^*$。
